\uuid{INyE}
\exo7id{3481}
\titre{exo7 3481}
\auteur{quercia}
\organisation{exo7}
\datecreate{2010-03-10}
\isIndication{false}
\isCorrection{false}
\chapitre{Déterminant, système linéaire}
\sousChapitre{Système linéaire, rang}
\cpge{mpsi-mp2i.calcul-algebrique-trigonometrie.b, pcsi.calcul-algebrique-trigonometrie.b, ptsi.calcul-algebrique-trigonometrie.b, mpsi-mp2i.calcul-matriciel-systemes.e, pcsi.determinants.c, ptsi.determinants.a, mpsi-mp2i.determinants.d, pcsi.calcul-matriciel-systemes-lineaires.d, ptsi.calcul-matriciel-systemes-lineaires.d, bcpst1.matrices}
\module{Algèbre}
\niveau{L2}
\difficulte{}

\contenu{
\texte{
Soit $M=(a_{ij}) \in \mathcal{M}_n(\C)$. On dit que $M$ est {\it à diagonale dominante\/}
si : $\forall\ i,\ |a_{ii}| > \sum_{j\ne i}|a_{ij}|$.
}
\begin{enumerate}
    \item \question{On transforme $M$ en
    $M' = \begin{pmatrix} a_{11} &\ a_{12}&\dots &a_{1n} \cr
                     0                         \cr
                   \vdots &        &M_1        \cr
                     0                         \cr \end{pmatrix}$
    par la méthode du pivot. Montrer que $M_1$ est à diagonale dominante.}
    \item \question{En déduire que $M$ est inversible.}
\end{enumerate}
}
